% Part: set-theory
% Chapter: ordinals
% Section: introduction

\documentclass[../../../include/open-logic-section]{subfiles}

\begin{document}

\olfileid{sth}{ordinals}{intro}

\olsection{પરિચય}

\olref[z][]{chap}માં આપણે \stagesinf{}ના સ્વરૂપે પદાનુક્રમનો
પહેલો અનંત તબક્કો હોય એમ સ્વીકાર્યું હતું (અનંતતાની સ્વયંસિદ્ધિ
પણ જુઓ). જોકે \stagessucc{} અપાયેલું હોય ત્યારે આપણે અનંત
તબક્કે અટકી શકતા નથી; આગળ વધતા રહેવું પડે. તેથી પહેલા
અનંત તબક્કા પછીના તબક્કે, પહેલા અનંત તબક્કે ઉપલબ્ધ થયેલા
ગણોના બધા સંભવિત સંગ્રહો રચીએ છીએ; અને ફરી એ જ કરીએ;
અને ફરી; અને ફરી; \dots

અહીં પરોક્ષ રીતે એવું થયું છે કે આપણે સંખ્યાનો એક
“સહજ” ખ્યાલ વાપરવાનું શરૂ કર્યું છે, જેમાં બધી પ્રાકૃતિક
સંખ્યાઓ \emph{પછી} પણ સંખ્યાઓ હોઈ શકે છે. ખાસ કરીને,
આ ખ્યાલ \emph{અનંતાતીત ક્રમવાચક સંખ્યાનો} છે.
આ વિચારને વધુ ચોક્કસ બનાવવો એ આ પ્રકરણનો હેતુ છે.
આપણે ક્રમવાચક સંખ્યાનો સામાન્ય ખ્યાલ તપાસીશું, અને પછી
અમુક ગણોને આપણી ક્રમવાચક સંખ્યાઓ તરીકે સ્પષ્ટ રીતે
વ્યાખ્યાયિત કરીશું.

\end{document}
